\documentclass{article}
\usepackage{graphicx} % Required for inserting images
\usepackage{amsmath}

\title{11-6 ECE544 Weekly Report}
\author{Gregory Stokes, Randy Peirce}
\date{November 6, 2026}

\begin{document}

\maketitle

\section{Progress and tasks completed this week:}
\subsection{Gregory Stokes}
    \begin{itemize}
        \item Attempted to convert sax netlists to GDS files
        \item Learned how to run LiDAR in the system
    \end{itemize}
\subsection{Randy Peirce}
    \begin{itemize}
        \item Established 2D baseline routing harness for Benes $4\times4$ switch matrix in gdsfactory
        \item Resolved component placement overlap using structured stage/pitch spacing ($400\,\mu\text{m} \times 250\,\mu\text{m}$)
        \item Configured dedicated environment and successfully ran native LiDAR 2D router on benchmark circuits
        \item Analyzed 2D layout in KLayout to identify waveguide crossing bottlenecks requiring 3D layer transitions
    \end{itemize}

\section{Plan for Next Week:}
\subsection{Gregory Stokes}
    \begin{itemize}
        \item Create a new virtual environment and run LiDAR on the system
        \item Understand how the code does DRC
        \item Understand the difference between DRC 2D vs 3D
        \item Create a DRC for a 3D system
    \end{itemize}

\subsection{Randy Peirce}
    \begin{itemize}
        \item Implement 3D loss-aware cost function in Python ($1550\,\text{nm}$ calibration for Si vs. SiN)
        \item Define 3D routing state space with vertical taper footprint and keep-out constraints
        \item Implement bounding-box search corridors and U-turn pruning heuristics to ensure scalability
        \item Run initial 3D single-net pathfinding on the Benes $4\times4$ crossover nets to demonstrate crossing elimination
    \end{itemize}

\end{document}
